% GENERATED FILE. Edit canonical lessons, then run npm run generate:books.
% canonical-source-hash: fnv1a64:582004ab397dc291
% canonical-lessons: TE-C367-sithilalu, TE-C367-puratana, TE-C367-tavvakalu, TE-C367-avasesalu, TE-C367-velikitiyu

\chapter{Digging Up the Past}
\label{ch:te-digging-up-the-past}
\hlchaptermodality{367}

\begin{chapteropening}
\textbf{By the end of this chapter:} \emph{I can say ruins, ancient, excavations, remains and to unearth.}

Complete the last lesson of chapter 367: I can say ruins, ancient, excavations, remains and to unearth.
\end{chapteropening}

\section[\texorpdfstring{\te{శిథిలాలు} (śithilālu)}{శిథిలాలు (śithilālu)}]{\te{శిథిలాలు} (śithilālu) — ruins}

\label{lesson:TE-C367-sithilalu}



\begin{quote}
Before the new one: say the Telugu for a lord, then the Telugu for to occupy.
\end{quote}

\subsection*{You'll want to know: \te{శిథిలాలు}}
\textbf{\te{శిథిలాలు}} — \emph{śithilālu} — \textquotedblleft{}ruins\textquotedblright{}.

\begin{quote}
\te{అడవిలో} \te{పాత} \te{కోట} \te{శిథిలాలు} \te{ఉన్నాయి}.
\end{quote}

\begin{quote}
\emph{aḍavilō pāta kōṭa śithilālu unnāyi.}
\end{quote}

\begin{quote}
There are ruins of an old fort in the forest.
\end{quote}

A plural noun; \te{శిథిల}, loose or decayed, makes the adjective.

\begin{cousinweb}
\textbf{\te{శిథిలాలు}} (\emph{śithilālu}) is built on Sanskrit \emph{śithila}, \textquotedblleft{}loose, slack\textquotedblright{} (Monier-Williams): what has come loose and fallen.
\end{cousinweb}

\begin{grammarlens}[title={What you've built}]
One more word for telling what happened.
\end{grammarlens}

\subsection*{Your turn}
\begin{itemize}
\raggedright
  \item \emph{Say it:} \emph{śithilālu}
  \item \emph{Say it:} \emph{śithilālu}, once more
  \item \emph{Say it:} \emph{aḍavilō pāta kōṭa śithilālu unnāyi.} — There are ruins of an old fort in the forest.
  \item \emph{Recall:} say the Telugu for a lord, then the Telugu for to occupy, then say \emph{śithilālu} again
\end{itemize}

\begin{tcolorbox}[breakable,colback=teal!4,colframe=teal!35!black,title={Before you move on}]
What does \te{శిథిలాలు} mean? (\textquotedblleft{}Ruins\textquotedblright{}.) Say it once more.
\end{tcolorbox}
\section[\texorpdfstring{\te{పురాతన} (purātana)}{పురాతన (purātana)}]{\te{పురాతన} (purātana) — ancient; very old}

\label{lesson:TE-C367-puratana}



\begin{quote}
Before the new one: say the Telugu for to occupy, then the Telugu for ruins.
\end{quote}

\subsection*{You'll want to know: \te{పురాతన}}
\textbf{\te{పురాతన}} — \emph{purātana} — \textquotedblleft{}ancient; very old\textquotedblright{}.

\begin{quote}
\te{అది} \te{ఒక} \te{పురాతన} \te{దేవాలయం}.
\end{quote}

\begin{quote}
\emph{adi oka purātana dēvālayaṁ.}
\end{quote}

\begin{quote}
That is an ancient temple.
\end{quote}

An adjective that goes before its noun: \te{పురాతన} \te{కట్టడం}, an ancient building.

\begin{cousinweb}
\textbf{\te{పురాతన}} (\emph{purātana}) is borrowed from Sanskrit \emph{purātana}, \textquotedblleft{}ancient\textquotedblright{} (CDIAL 8285).
\end{cousinweb}

\begin{grammarlens}[title={What you've built}]
One more word for telling what happened.
\end{grammarlens}

\subsection*{Your turn}
\begin{itemize}
\raggedright
  \item \emph{Say it:} \emph{purātana}
  \item \emph{Say it:} \emph{purātana}, once more
  \item \emph{Say it:} \emph{adi oka purātana dēvālayaṁ.} — That is an ancient temple.
  \item \emph{Recall:} say the Telugu for to occupy, then the Telugu for ruins, then say \emph{purātana} again
\end{itemize}

\begin{tcolorbox}[breakable,colback=teal!4,colframe=teal!35!black,title={Before you move on}]
What does \te{పురాతన} mean? (\textquotedblleft{}Ancient; very old\textquotedblright{}.) Say it once more.
\end{tcolorbox}
\section[\texorpdfstring{\te{తవ్వకాలు} (tavvakālu)}{తవ్వకాలు (tavvakālu)}]{\te{తవ్వకాలు} (tavvakālu) — excavations}

\label{lesson:TE-C367-tavvakalu}



\begin{quote}
Before the new one: say the Telugu for ruins, then the Telugu for ancient.
\end{quote}

\subsection*{You'll want to know: \te{తవ్వకాలు}}
\textbf{\te{తవ్వకాలు}} — \emph{tavvakālu} — \textquotedblleft{}excavations\textquotedblright{}.

\begin{quote}
\te{తవ్వకాలలో} \te{పాత} \te{నాణేలు} \te{దొరికాయి}.
\end{quote}

\begin{quote}
\emph{tavvakālalō pāta nāṇēlu dorikāyi.}
\end{quote}

\begin{quote}
Old coins were found in the excavations.
\end{quote}

From \te{తవ్వు}, to dig.

\begin{cousinweb}
\textbf{\te{తవ్వకాలు}} (\emph{tavvakālu}) is native Dravidian, the plural noun of \emph{tavvu}, \textquotedblleft{}to dig, excavate\textquotedblright{} (DEDR 3122).
\end{cousinweb}

\begin{grammarlens}[title={What you've built}]
One more word for telling what happened.
\end{grammarlens}

\subsection*{Your turn}
\begin{itemize}
\raggedright
  \item \emph{Say it:} \emph{tavvakālu}
  \item \emph{Say it:} \emph{tavvakālu}, once more
  \item \emph{Say it:} \emph{tavvakālalō pāta nāṇēlu dorikāyi.} — Old coins were found in the excavations.
  \item \emph{Recall:} say the Telugu for ruins, then the Telugu for ancient, then say \emph{tavvakālu} again
\end{itemize}

\begin{tcolorbox}[breakable,colback=teal!4,colframe=teal!35!black,title={Before you move on}]
What does \te{తవ్వకాలు} mean? (\textquotedblleft{}Excavations\textquotedblright{}.) Say it once more.
\end{tcolorbox}
\section[\texorpdfstring{\te{అవశేషాలు} (avaśēṣālu)}{అవశేషాలు (avaśēṣālu)}]{\te{అవశేషాలు} (avaśēṣālu) — remains; relics}

\label{lesson:TE-C367-avasesalu}



\begin{quote}
Before the new one: say the Telugu for ancient, then the Telugu for excavations.
\end{quote}

\subsection*{You'll want to know: \te{అవశేషాలు}}
\textbf{\te{అవశేషాలు}} — \emph{avaśēṣālu} — \textquotedblleft{}remains; relics\textquotedblright{}.

\begin{quote}
\te{అక్కడ} \te{పాత} \te{నగరం} \te{అవశేషాలు} \te{దొరికాయి}.
\end{quote}

\begin{quote}
\emph{akkaḍa pāta nagaraṁ avaśēṣālu dorikāyi.}
\end{quote}

\begin{quote}
The remains of an old city were found there.
\end{quote}

\begin{cousinweb}
\textbf{\te{అవశేషాలు}} (\emph{avaśēṣālu}) is built on Sanskrit \emph{avaśeṣa}, \textquotedblleft{}leavings, remainder\textquotedblright{} (Monier-Williams).
\end{cousinweb}

\begin{grammarlens}[title={What you've built}]
One more word for telling what happened.
\end{grammarlens}

\subsection*{Your turn}
\begin{itemize}
\raggedright
  \item \emph{Say it:} \emph{avaśēṣālu}
  \item \emph{Say it:} \emph{avaśēṣālu}, once more
  \item \emph{Say it:} \emph{akkaḍa pāta nagaraṁ avaśēṣālu dorikāyi.} — The remains of an old city were found there.
  \item \emph{Recall:} say the Telugu for ancient, then the Telugu for excavations, then say \emph{avaśēṣālu} again
\end{itemize}

\begin{tcolorbox}[breakable,colback=teal!4,colframe=teal!35!black,title={Before you move on}]
What does \te{అవశేషాలు} mean? (\textquotedblleft{}Remains; relics\textquotedblright{}.) Say it once more.
\end{tcolorbox}
\section[\texorpdfstring{\te{వెలికితీయు} (velikitīyu)}{వెలికితీయు (velikitīyu)}]{\te{వెలికితీయు} (velikitīyu) — to unearth; to dig out}

\label{lesson:TE-C367-velikitiyu}



\begin{quote}
Before the new one: say the Telugu for excavations, then the Telugu for remains.
\end{quote}

\subsection*{You'll want to know: \te{వెలికితీయు}}
\textbf{\te{వెలికితీయు}} — \emph{velikitīyu} — \textquotedblleft{}to unearth; to dig out\textquotedblright{}.

\begin{quote}
\te{శాస్త్రవేత్తలు} \te{ఒక} \te{పాత} \te{కుండను} \te{వెలికితీశారు}.
\end{quote}

\begin{quote}
\emph{śāstravēttalu oka pāta kuṇḍanu velikitīśāru.}
\end{quote}

\begin{quote}
The scientists unearthed an old pot.
\end{quote}

Past (I …): \te{వెలికితీశాను}, \emph{velikitīśānu}.

Literally to pull to the outside; also to bring a secret to light.

\begin{cousinweb}
\textbf{\te{వెలికితీయు}} (\emph{velikitīyu}) is a native compound: \emph{veliki}, \textquotedblleft{}to the outside\textquotedblright{}, from \emph{veli}, \textquotedblleft{}the outside\textquotedblright{} (DEDR 5498), and \emph{tīyu}, \textquotedblleft{}to draw, pull out\textquotedblright{}.
\end{cousinweb}

\begin{grammarlens}[title={What you've built}]
One more word for telling what happened.
\end{grammarlens}

\subsection*{Your turn}
\begin{itemize}
\raggedright
  \item \emph{Say it:} \emph{velikitīyu}
  \item \emph{Say it:} \emph{velikitīyu}, once more
  \item \emph{Say it:} \emph{śāstravēttalu oka pāta kuṇḍanu velikitīśāru.} — The scientists unearthed an old pot.
  \item \emph{Recall:} say the Telugu for excavations, then the Telugu for remains, then say \emph{velikitīyu} again
\end{itemize}

\begin{tcolorbox}[breakable,colback=teal!4,colframe=teal!35!black,title={Before you move on}]
What does \te{వెలికితీయు} mean? (\textquotedblleft{}To unearth; to dig out\textquotedblright{}.) Say it once more.
\end{tcolorbox}
